\chapter{Making decisions about unplanned pregnancies}
\index{Making decisions about unplanned pregnancies}

\section*{Definition and Overview}
Making decisions about unplanned pregnancies is a critical, complex clinical and personal process that occurs when a pregnancy is unexpected, unintended, or becomes unsustainable due to changing life circumstances. Clinically, this requires a supportive, non-judgemental, and evidence-based framework to assist individuals in exploring their options. The three primary pathways available are continuing the pregnancy to parent, continuing the pregnancy to pursue adoption or alternative care arrangements (such as permanent care or foster care), and terminating the pregnancy via medical or surgical abortion. Access to unbiased, nondirective counselling, accurate reproductive health information, and timely medical care is essential, as delays in decision-making can limit clinical options, particularly regarding the methods and availability of termination or adoption pathways.

\section*{Epidemiology}
Unplanned pregnancy is a common global health phenomenon. Approximately 40\% to 50\% of all pregnancies worldwide are estimated to be unintended, and a significant proportion occur despite the active use of contraception. In many countries, it is estimated that one in four women will experience an unplanned pregnancy at some point in their reproductive life. Rates of unplanned pregnancy vary across demographics, with higher incidences observed among younger women, those residing in rural or remote areas, and socioeconomically disadvantaged populations. The decision-making outcomes also vary; approximately half of unplanned pregnancies in many countries are estimated to result in termination, while the remainder are carried to term for parenting or, less frequently, adoption.

\section*{Aetiology and Pathophysiology}
Unplanned pregnancies arise from various biological, behavioural, and social factors that lead to fertilisation in the absence of a planned intent to conceive. The clinical and pathophysiological contexts that lead to decision-making include:
\begin{itemize}
    \item \textbf{Contraceptive failure:} The incorrect, inconsistent, or unexpected failure of contraceptive methods, which occurs with varying rates depending on the method's Pearl Index (such as condom breakage, missed oral contraceptive pills, or displacement of an intrauterine device).
    \item \textbf{Non-use of contraception:} Lack of contraceptive coverage due to lack of access, educational barriers, side effects, or unexpected intercourse.
    \item \textbf{Reproductive coercion:} Control over a partner's reproductive choices by an abusive partner, including forced intercourse, tampering with birth control, or pressure to carry a pregnancy to term.
    \item \textbf{Changes in maternal or fetal health:} A planned pregnancy that becomes unplanned or complicated due to the acute diagnosis of severe maternal illness, fetal abnormalities, or genetic disorders during gestation.
\end{itemize}

\section*{Clinical Features}
Individuals presenting with an unplanned pregnancy exhibit a range of clinical and psychosocial features that must be addressed during consultation. The presentation includes:
\begin{itemize}
    \item \textbf{Amenorrhoea and early pregnancy symptoms:} Absence of menstruation, nausea, breast tenderness, fatigue, and frequent urination, which prompt pregnancy testing.
    \item \textbf{Psychological distress:} Feelings of shock, anxiety, denial, grief, confusion, or ambivalence regarding the future, which require immediate psychological support.
    \item \textbf{Socioeconomic concerns:} Anxiety regarding financial stability, housing, relationship status, or the impact on education and career pathways.
    \item \textbf{Urgency and time constraints:} Awareness of gestational limits for medical or surgical interventions, creating a sense of pressure and cognitive overload.
\end{itemize}

\section*{Differential Diagnosis}
In a clinical setting, it is important to confirm intrauterine pregnancy and evaluate the patient's capacity and safety. The differential diagnoses and safety concerns to address include:
\begin{itemize}
    \item \textbf{Ectopic pregnancy:} A life-threatening extrauterine pregnancy (commonly tubal) that must be ruled out in patients presenting with pelvic pain or spotting, as it requires immediate medical or surgical intervention rather than standard termination or antenatal care.
    \item \textbf{Molar pregnancy:} A gestational trophoblastic disease that presents with elevated beta-hCG levels, severe nausea, and vaginal bleeding, necessitating surgical evacuation.
    \item \textbf{Severe mental health disorder:} Acute exacerbation of psychiatric illness that may impair decision-making capacity, requiring urgent psychiatric assessment and support.
    \item \textbf{Domestic and family violence:} Coercive control or physical abuse that influences the patient's decision or compromises their safety, requiring immediate social work or protective intervention.
\end{itemize}

\section*{When to Seek Medical Care}
Individuals should seek prompt medical consultation upon discovering an unplanned pregnancy to confirm gestational age and discuss options.

\textbf{Contact your local emergency services for:}
\begin{itemize}
    \item \textbf{Severe abdominal or pelvic pain:} Sharp, localised, or worsening pain, particularly on one side, which may indicate an ectopic pregnancy rupture.
    \item \textbf{Heavy vaginal bleeding:} Severe haemorrhage with or without clots, potentially indicating a miscarriage, ectopic pregnancy, or other uterine bleeding.
    \item \textbf{Sepsis:} High fever, chills, rapid heart rate, or feeling extremely unwell, which could point to a pelvic infection.
    \item \textbf{Signs of shock:} Dizziness, fainting, pale skin, cold sweats, or rapid breathing associated with internal bleeding.
\end{itemize}

\section*{Diagnosis and Investigations}
Accurate diagnosis and gestational dating are crucial to guide clinical decisions. Investigations include:
\begin{itemize}
    \item \textbf{Urine and serum beta-hCG testing:} To confirm pregnancy and monitor hormone levels if ectopic pregnancy or miscarriage is suspected.
    \item \textbf{Pelvic ultrasound:} The definitive investigation to confirm intrauterine location, determine accurate gestational age, and assess viability.
    \item \textbf{Blood group and antibody screening:} Essential to determine Rhesus (Rh) status; Rh-negative patients require anti-D immunoglobulin if undergoing termination or experiencing bleeding.
    \item \textbf{Screening for sexually transmitted infections (STIs):} Routine testing for chlamydia, gonorrhoea, and blood-borne viruses, particularly prior to surgical termination, to reduce the risk of post-procedure pelvic inflammatory disease.
    \item \textbf{Psychosocial assessment:} Evaluation of support systems, safety, coercion, and mental health status to ensure informed, voluntary decision-making.
\end{itemize}

\section*{Management}
The management of unplanned pregnancy involves providing comprehensive, objective information, counselling, and clinical pathways for the chosen option. Management strategies include:
\begin{itemize}
    \item \textbf{Nondirective pregnancy options counselling:} Structured, unbiased counselling that supports the patient to explore their feelings and values regarding parenting, adoption, and termination without pressure or judgement.
    \item \textbf{Termination of pregnancy pathways:} Facilitating referral to medical abortion (using mifepristone and misoprostol) or surgical abortion (suction curettage or dilation and evacuation), depending on gestational age and patient choice.
    \item \textbf{Antenatal care pathway:} Setting up early, comprehensive prenatal care for patients choosing to continue the pregnancy to parent.
    \item \textbf{Adoption and foster care pathways:} Referral to specialised social workers and organisations to discuss open, semi-open, or closed adoption, or temporary foster care arrangements.
    \item \textbf{Contraceptive counselling:} Detailed discussion and provision of reliable contraception, especially long-acting reversible contraceptives (LARCs) like implants or intrauterine devices, to prevent future unintended pregnancies.
\end{itemize}

\section*{Complications}
The complications associated with the decision-making process are largely psychosocial, but delayed choice can lead to clinical risks. These include:
\begin{itemize}
    \item \textbf{Delayed access to care:} Late-term presentation that limits the availability of medical termination, requiring more invasive surgical options or travel to specialised facilities.
    \item \textbf{Psychological distress:} Long-term anxiety, depression, or post-traumatic stress, particularly if the decision was coerced, rushed, or lacked adequate support.
    \item \textbf{Relationship breakdown:} Conflict or separation due to divergent views between partners regarding the pregnancy.
    \item \textbf{Stigmatisation and isolation:} Social isolation or shame arising from cultural, religious, or familial disapproval of termination or adoption.
\end{itemize}

\section*{Prognosis and Outcomes}
The majority of individuals who make decisions about unplanned pregnancies adapt well and do not experience long-term adverse psychological effects, provided the decision was made autonomously, with adequate support, and without coercion. Post-decision outcomes depend on the pathway chosen. Those who parent benefit from early engagement with maternal health services, while those who undergo termination generally report relief and a rapid return to physical health. Continued access to mental health support, peer counselling, and post-abortion care is crucial for the minority of patients who experience complex grief or adjustment disorders.

\section*{Prevention and Self-Care}
Preventative strategies focus on reducing the risk of unplanned pregnancy and supporting emotional well-being during decision-making. Recommendations include:
\begin{itemize}
    \item \textbf{Contraceptive adherence:} Understanding and consistently using highly effective contraceptive methods, particularly long-acting options like implants or IUDs.
    \item \textbf{Emergency contraception:} Accessing emergency contraceptive pills (such as levonorgestrel or ulipristal acetate) or emergency copper IUD insertion within 72 to 120 hours of unprotected intercourse.
    \item \textbf{Accessing support networks:} Reaching out to trusted friends, family, or professional counsellors to discuss feelings and options in a safe environment.
    \item \textbf{Self-education:} Obtaining factual, non-biased information from reputable medical websites and family planning organisations.
    \item \textbf{Prioritising physical safety:} Leaving unsafe relationships and seeking refuge through domestic violence support services if reproductive coercion or abuse is present.
\end{itemize}

\section*{Sources and Further Reading}
The following reputable organisations provide unbiased information, support, and resources on unplanned pregnancies:
\begin{itemize}
    \item Healthdirect Australia. \textit{Support options and decision-making resources for unplanned pregnancy.} https://www.healthdirect.gov.au/
    \item MSI Reproductive Choices Australia. \textit{Support services, counselling, and termination information.} https://www.msichoices.org.au/
    \item Jean Hailes for Women's Health. \textit{Information on contraception, sexual health, and unplanned pregnancy.} https://www.jeanhailes.org.au/
    \item NHS. \textit{Unbiased guide to pregnancy options, abortion, and adoption.} https://www.nhs.uk/
    \item 1800MyOptions (Victoria). \textit{Information and phone support for contraception, pregnancy options, and services.} https://www.1800myoptions.org.au/
\end{itemize}
